\documentclass[10pt,a4paper]{article}
\usepackage[margin=23mm]{geometry}
\usepackage{amsmath,amssymb,booktabs,array,longtable}
\usepackage[hidelinks]{hyperref}
\setlength{\parindent}{0pt}
\setlength{\parskip}{4pt}
\setlength{\emergencystretch}{2em}
\newcommand{\dd}{\mathrm d}
\newcommand{\pos}[1]{[#1]_+}
\title{Urban service allocation:\\traffic and sensitivity specification}
\author{Urban traffic methods}
\date{28 October 2026}
\begin{document}
\maketitle

\textbf{Purpose and scope.}
This document specifies a discrete continuous-mass model for an urban
signal-control task. It defines lane admission, order-dependent access,
service and complete population-time sensitivity. The representative
approach has two lanes upstream of a finite left-turn bay and three lanes
at the stop line. Numerical calibration, baseline implementation checks
and empirical validation are separate from this mathematical definition.

\section{Traffic world, action and physical ownership}

Retain one isolated four-leg junction, fixed routes and four protected
effective-service groups: NS through/right, NS left, EW through/right,
EW left. With physical capacity $C_\ell$ and saturation $s_{\ell m}$,
\begin{equation}
 g_p\geq0,\quad \sum_{p=1}^4g_p=1,\qquad
 \mu_{\ell m}(g)=\min\{C_\ell,s_{\ell m}g_{p(m)}\}.
 \label{eq:action}
\end{equation}
Mass is continuous vehicle-equivalent quantity; $C,s,\mu$ are veh-eq/s.
A feasible direction preserves the simplex. Route labels, permissions,
topology and source-order skeleton are structural. No phase waveform,
integerization, microscopic lane changing or endogenous route choice is
introduced. Control uses constant $g$ during the declared control window,
then a common candidate-independent clearance allocation.

The representative approach has two parallel feeder lanes and three
stop-line lanes:
\begin{equation}
 U\longrightarrow\{F_1,F_2\},\qquad
 F_1\longrightarrow\{B,T_1\},\qquad F_2\longrightarrow T_2 .
 \label{eq:geometry}
\end{equation}
$U$ is the aggregate upstream road. Within the local zone, $F_1$ carries
left/through traffic and $F_2$ through/right traffic. The internal bay
opening divides $F_1$ into left-only bay $B$ and through-only continuation
$T_1$; the same cross-section divides $F_2$ from shared through/right
continuation $T_2$. In the symmetric core
$L_{F_1}=L_{F_2}$ and $L_B=L_{T_1}=L_{T_2}$, giving equal total approach
path length. All sections own disjoint physical road space and storage.

Left traffic may enter only $F_1$, then $B$; right traffic only $F_2$,
then $T_2$. Through traffic may choose $F_1$ or $F_2$ at actual admission
from $U$, but follows its chosen lane thereafter: $F_1$ to $T_1$ or
$F_2$ to $T_2$. There are no cross-feeder lane-change edges and no direct
$U$ to $B,T_1,T_2$ edges. The queue estimate at the choice boundary uses
current receiving-feeder mass and earlier actual admissions, as in the
allocation rule below; no downstream look-ahead or new driver cost is added.
The two physical lanes can therefore differ in accessibility without
allowing traffic to escape a lane it already occupies.

V0 uses movement-exclusive full-length paths. V1 uses a shared lane.
V2/V3 share the representative bay geometry and permissions: V2 suppresses
order-sensitive obstruction by construction, whereas V3 tests it.
V4 supplies predeclared held-out variation. The smaller single-feeder
case remains a local mechanism diagnostic. Permissions may differ
between scenario families, but not between compared models or paired
order cases of the same scenario.

\textbf{Processing order and physical lane order.}
An unassigned link-end or external-source pool has an ordered processing
record. A realized local lane has an actual first-in, first-out sequence.
The input record governs allocation inspections; the occupied lane also
imposes front blocking. A waiting head in $F_1$ cannot be passed by a
later resident, while an independent $F_2$ has its own available path.
A record cannot be relabeled as unassigned to permit physical overtaking.

\section{Propagation and the step clock}

Each physical section retains cumulative counts $N^+_{e,k},N^-_{e,k}$
and occupancy $n_{e,k}=N^+_{e,k}-N^-_{e,k}$ at $t_k=k\Delta t$.
Use triangular-FD capacity $C_e$, free-flow delay $\tau^f_e$, backward
delay $\tau^b_e$ and storage $K_e=C_e(\tau^f_e+\tau^b_e)$, with
$\Delta t\leq\min_e(\tau^f_e,\tau^b_e)$.
Road capacities, delays and legal saturation rates are finite and positive.
An incompatible short section requires a compatible discretization;
no fictitious delay, clipping or smoothing is inserted.

With empty prehistory and $\theta=(t-t_h)/\Delta t$,
\begin{equation}
 {\cal H}^{\pm}_e(t)
 =(1-\theta)N^\pm_{e,h}+\theta N^\pm_{e,h+1},
 \qquad t_h\leq t\leq t_{h+1}.
 \label{eq:history}
\end{equation}
The sending count subtracts prior departures from entries early enough
to reach the exit. The receiving count starts with storage, adds departures
whose space-release wave has reached the entrance, and subtracts all
previous entries. Both counts are bounded by capacity times timestep.
The inherited one-step limits are
\begin{align}
 A_{e,k}&=\pos{{\cal H}^+_e(t_{k+1}-\tau^f_e)-N^-_{e,k}},
 &S_{e,k}&=\min(C_e\Delta t,A_{e,k}),\label{eq:sending}\\
 R_{e,k}&=\min\{C_e\Delta t,\,
 \pos{K_e+{\cal H}^-_e(t_{k+1}-\tau^b_e)-N^+_{e,k}}\}.
 \label{eq:receiving}
\end{align}
The positive part sets negative availability to zero. Free-flow delay
controls downstream travel; backward delay controls how soon new room
is available upstream. $A$ is ready resident mass for this update;
$S,R$ are mass allowances,
not instantaneous rates. Queries use already stored boundaries.
New inflow cannot traverse its receiving section in the same update.

This is a timestep model. The internal head-service calculation produces
endpoint totals; it does not add propagation-history knots.
Readiness uses the inherited discrete sending convention, rather than
claiming every selected resident was ready at every instant in the step.
No event-history refinement is adopted.

These equations and the left-interval objective are inherited from the
approved repository M0 transcription, checked against the local LTM and
objective code. Their application to the internal section graph is the
declared urban composition, not existing implemented M2 behavior.

\section{Ledger, eligibility and lane membership}

Store positive-mass cohorts, removing exact zero entries without a cutoff.
A cohort stores route-table reference, route position, continuous mass,
current owner and affine eligibility provenance. Current movement is
projected from its route. Internal $U/F_1/F_2/B/T_1/T_2$ crossings preserve route
position; the actual junction crossing advances it. Future movement
sequences are not duplicated in grid entries.

Only actual admission establishes receiving-lane membership. Unadmitted
mass stays with its old owner without a binding future-lane vector.
Traffic already in $F_1$ or $F_2$ retains that realized lane; its movement
and the declared connectivity determine its continuation. Resident traffic
cannot be reassigned across feeders or terminal lanes to improve balance.

Retain the exact source order of the unserved record. Before evaluating
an operator, coalesce exactly compatible adjacent representation fragments:
same route/position, owner, permissions and contiguous affine eligibility
with the M1 exact line condition. No tolerance-based mixing or route
merging is used. Serialization changes must preserve predictions; they
do not define a new physical traffic unit.

The movement grid is a view of the ordered ledger:
\begin{equation}
 M_{cm}=\sum_{r:m(r)=m}|I_r\cap B_c|,
 \label{eq:grid}
\end{equation}
where $I_r,B_c$ are intervals in cumulative resident mass. Rank is not
metres; mixed composition does not reconstruct exact head order.
Cells cover the full resident mass-rank range without truncating the
ledger. Storage, travel and spillback come from physical histories.

\textbf{Sending limits transfers, not candidate inspections.}
Form candidates from ready mass $A$, cutting an eligibility segment
exactly if needed. For an unassigned pool, scan these candidates and
spend at most $S$ in total. The candidate list is not restricted to the first $S$ units before access
is checked: an earlier blocked movement must not hide later ready traffic
with an independently accessible path. For a physical lane, the access rule below forces
actual served mass to be a prefix.

This is an explicit adaptation of M1's prefix extraction. After a
selection that skips an independently blocked movement, unserved mass
keeps its original entry/travel history and relative residual order.
Waiting for admission does not restart its travel delay. Every past departure was eligible when served, and eligibility
is monotone. Total eligible residual mass therefore remains
${\cal H}^+(t_{k+1}-\tau^f)-N^-_k$ even if eligible departures were not
a global processing-order prefix. The adaptation is not claimed to exist
in current M1 code, nor to establish passing on a represented single lane.

\section{Allocation, unassigned residuals and physical access}

Let $a$ be the sending owner, $r$ a ready cohort of mass $q_r$ and
movement $m_r$, and $\ell$ a receiving lane/section.
Permission $P_{m\ell}\in\{0,1\}$ and reachability
$E_{r\ell}\in\{0,1\}$ encode actual paths/access, not a future green plan.

Initialize residual sending $s=S_{a,k}$, residual receiving
$R^{\rm rem}_\ell=R_{\ell,k}$ and perceived loads $h_\ell=n_{\ell,k}$.
Process cohorts in recorded order. Let ${\cal B}$ contain movements
already left partly unserved in this scan. Each sending owner has its own
scan and blocked-movement set; one inbound leg does not block another.
If $m_r\in{\cal B}$ set $V_r=x_{r\ell}=0$; otherwise admit the largest
feasible total $V_r$, limited by packet mass, remaining sending and
legal reachable receiving. Distribute that total to reduce differences
in current lane loads, using a common filling level $\lambda_r$:

\begin{align}
 u_{r\ell}&=P_{m_r\ell}E_{r\ell}R^{\rm rem}_\ell,\qquad
 V_r=\min\{q_r,s,\sum_\ell u_{r\ell}\},\nonumber\\
 x_r&=\underset{0\leq x_{r\ell}\leq u_{r\ell},\
                    \sum_\ell x_{r\ell}=V_r}{\arg\min}
          \frac12\sum_{\ell:P_{m_r\ell}=1}(h_\ell+x_{r\ell})^2,
 \nonumber\\
 x_{r\ell}&=\min\{u_{r\ell},\pos{\lambda_r-h_\ell}\},\qquad
                    \sum_\ell x_{r\ell}=V_r .
 \label{eq:allocation}
\end{align}
The solver chooses $\lambda_r$ so the lane additions sum to $V_r$.
Its units are vehicle-equivalent mass: $\lambda_r-h_\ell$ is the mass
needed to bring a lane to that load. The positive part excludes lanes
already above the level and $u_{r\ell}$ caps each feasible admission.
This is a mass-balancing rule, not a stop-line service cap. Illegal lanes
have zero addition. Strict convexity selects unique additions; the filling
level itself need not be unique when all usable lanes are capped.
For $V=0$, use $x=0$. Commit only actual entry:
\begin{equation}
 s\leftarrow s-V_r,\quad
 R^{\rm rem}_\ell\leftarrow R^{\rm rem}_\ell-x_{r\ell},\quad
 h_\ell\leftarrow h_\ell+x_{r\ell},\quad
 q_r^{\rm rem}=q_r-V_r .
 \label{eq:allocation-update}
\end{equation}
If $q_r^{\rm rem}>0$, retain it at its original position and add $m_r$
to ${\cal B}$ for this update. Later cohorts of that movement wait;
other movements are considered. Reset ${\cal B}$ next update.
The retained record then puts the older residual first within its
movement. There is no waiting-target allocation or old-target rationing law.

$h$ is the selected queue-perception estimate: pre-state local lane mass
plus earlier actual admissions. It contains no invented waiting targets
and does not subtract forecast signal discharge. Actual next occupancy
separately includes that step's exits. Raw queue mass is the first
perception model; no service-aware cost or unequal preference is added.

\subsection*{Access in the admitted layouts}

For an unassigned pool with independently accessible receiving lanes,
$E_{r\ell}=1$ on those paths. Movement labels alone do not create access.
This applies at $U\to\{F_1,F_2\}$ on a packet's legal feeder choices,
and to a source feeding independent V0 paths. Once mass occupies $F_1$
or $F_2$, it cannot use the other feeder's capacity.

For an occupied single lane, including $F_1$ at its fork,
\begin{equation}
 E_{r\ell}
 =\mathbf1\{\text{every earlier resident cohort has fully crossed}\}
 \quad\text{on a legal outgoing edge}.
 \label{eq:physical-access}
\end{equation}
An older physical residual prevents later physical entry even while the
allocator inspects later records. Several actual lanes have separate
ledgers/access conditions, not one global FIFO.

Movement-specific blocking governs the unassigned scan; actual lane
occupation adds the physical access constraint. At $U$, a full $F_1$
prevents left entry while through traffic may still use reachable $F_2$.
At the fork, a through head in $F_1$ can obstruct a later turner despite
an empty bay. A full bay can similarly leave a left head blocking later
through traffic in $F_1$. Independently resident $F_2$ traffic may continue
toward $T_2$ if its own sending, receiving and order permit it.

A zero left-turn stop-line rate is not automatically zero lane-entry room.
A left turner may join a reachable feeder or bay with room and subsequently
contribute to spillback. In this geometry it cannot turn left from $T_1$
or use $F_2,T_2$; a full bay has no alternative left-turn stop-line lane.
Each constraint refers to its actual physical boundary.

\subsection*{Direct definition checks}

The first case below is a local operator diagnostic, not a new primary
scenario or a claim that every aggregate road offers independent access.
For an independent-access pool with order $[L(0.4),T(0.6)]$, $A=1$,
$S=0.6$ and exclusive-destination receiving $(0,0.5)$, left entry is
zero and through entry $0.5$. Residual order is $[L(0.4),T(0.1)]$.
Truncating candidates to the first $S$ mass would expose only $0.2$
through mass and give the wrong result for this pool rule.

For actual $F_1$ take $[T(0.4),L(0.6)]$, $S_{F_1}=1$, $R_{T_1}=0.2$,
$R_B=1$. Through entry is $0.2$, left entry zero.
Reverse the order with the same totals and branch loads
$(n_B,n_{T_1})=(0,1.8)$. Left uses $B$ first, then through enters $0.2$:
total $0.8$ rather than $0.2$.
These branch bounds are compatible with
$\Delta t=\tau^f=\tau^b=1$, $C_B=C_{T_1}=1$, $K_B=K_{T_1}=2$.
If an independent ready $0.3$ in $F_2$ has enough receiving in $T_2$,
it enters in either case. Cross-section totals are then $0.5$ and $1.1$.
These are prepared hand states, not a direct initializer or rollout.

At the $U$ choice boundary, a compatible balancing case has feeder loads
$(1.2,1.0)$, receiving $(0.8,1.0)$, ready through mass $1$ and sending $0.6$.
Additions $(0.2,0.4)$ leave perceived loads $(1.4,1.4)$;
the remaining upstream $0.4$ is unassigned.

\section{Elapsed service and selected receiving rates}

A stop-line lane exposes its current ready resident head, with sending
quota $S_\ell$ and cap $c_\ell=\mu_{\ell m}(g)$.
The selected coarse receiving convention is
\begin{equation}
 \rho_j=R_{j,k}/\Delta t .
 \label{eq:receiving-rate}
\end{equation}
This is a timestep computational average, not an inferred within-step
physical supply curve.

For each output $j$, collect current heads by inbound leg $i$.
A leg-to-output movement here has one service group:
\begin{align}
 C_{ij}&=\sum_{\ell\in{\cal L}_{ij}}c_\ell,\quad
 w_{ij}=g_{p(i,j)},\quad W_j=\min(\rho_j,\sum_iC_{ij}),\nonumber\\
 r_{ij}&=\min(C_{ij},\theta_jw_{ij}),\quad
                    \sum_i r_{ij}=W_j,\label{eq:leg-sharing}\\
 r_\ell&=\min(c_\ell,\eta_{ij}),\quad
                    \sum_{\ell\in{\cal L}_{ij}}r_\ell=r_{ij}.
 \label{eq:lane-sharing}
\end{align}
Remove zero-cap competitors; zero $W$ gives zero rates.
Each leg is weighted once; its ready lanes share equally with caps/reuse.
Splitting a leg into more software lanes cannot multiply its leg priority.

Hold rates until the earliest head/quota exhaustion or step end,
advancing all lanes together. For residual head mass $q_\ell$,
\begin{equation}
 \delta=\min\left\{\Delta t-\sigma,\
 \min_{\ell:r_\ell>0}
       \frac{\min(q_\ell,S_\ell^{\rm rem})}{r_\ell}\right\}.
 \label{eq:head-time}
\end{equation}
Transfer $r_\ell\delta$, process tied events together, expose new heads
and recompute sharing. Stop when rates vanish or time ends.
A zero-rate head blocks its lane. Newly admitted mass is outside this
resident service tranche.

A later ready head joins current competition; earlier discharge is not
compensated to equalize final totals. Constant receiving removes a
separate consumable-pool exhaustion rule, but head clearance and changing
competitors remain sequential. Receiving reuse does not transfer one
signal group's service to another.

For V1's separate shared left/through service diagnostic, ready left
mass $q_L$ precedes sufficient through mass. With nonbinding receiving
and physical caps, its clearance time determines the remaining through
service time:
\begin{equation}
 y_L=q_L,\qquad
 y_T=s_Tg_T\left(\Delta t-\frac{q_L}{s_Lg_L}\right),
 \quad 0<\frac{q_L}{s_Lg_L}<\Delta t.
 \label{eq:head-example}
\end{equation}
With unit timestep/saturations, $q_L=0.2$ and $(g_L,g_T)=(0.4,0.6)$,
through discharge is $0.3$. A small transfer from through to left gives
through-discharge derivative $0.25$ on this branch: faster head removal
outweighs the smaller through rate. This is not a complete-TST benefit claim.

\section{Paired transfers, histories and feasibility}

A one-in/one-out internal edge transfers its ready prefix up to
$\min(S_{\rm up},R_{\rm down})$. Forks use allocation/access; the central
junction uses elapsed service. An external source adds $D_{o,k}$ to its
waiting record before admission, with budget $Q_{o,k}+D_{o,k}$.
A sink removes at most sending and its declared allowance.

Each internal entry interface has one sending owner; $U$ has two
receiving feeder choices sharing its total sending allowance. Shared receiving at the
central signal node uses the simultaneous current-head sharing law, rather
than an arbitrary serial order of inbound legs.
All nodes read pre-state histories/ledgers. Each section's sending and
receiving is spent once at its boundary. New admissions are appended
after selecting resident service; today's exits create no immediate
receiving credit. The cumulative update is
\begin{equation}
 N^+_{e,k+1}=N^+_{e,k}+b_{e,k},\quad
 N^-_{e,k+1}=N^-_{e,k}+y_{e,k},\quad
 n_{e,k+1}=n_{e,k}+b_{e,k}-y_{e,k}.
 \label{eq:state-update}
\end{equation}
Internal flow leaves one owner and enters the next with identical mass
and route identity. Source admission removes source waiting; sink exit
removes road mass. No residual is duplicated as lane stock.

For a receiving lane with total entrance $b>0$, append children in their
actual admission subsequence. A child at entrance rank $[r,r+q']$ gets
\begin{equation}
 \left[t_k+\Delta t\,\frac r b,\
       t_k+\Delta t\,\frac{r+q'}b\right],
 \qquad\text{eligibility shifted by }\tau^f_e .
 \label{eq:entry-time}
\end{equation}
For zero inflow create no segment. This is M1-style affine boundary
reconstruction applied to the new selection rule, without service-event
history knots. Serialize cross-input outgoing records by internal
service-interval start, then fixed input-leg index, lane index and source
cohort index. This orders macro bookkeeping, not rate grants or a new
physical zipper rule. Those roads lead to sinks here; their scalar
propagation/total time cannot acquire a later signal-order advantage
from that serialization.

The timestep domain gives $S_{e,k}\leq n_{e,k}$ and
$R_{e,k}\leq K_e-n_{e,k}$. Hence
\begin{equation}
 0\leq b_{e,k}\leq R_{e,k},\quad0\leq y_{e,k}\leq S_{e,k}
 \ \Longrightarrow\
 0\leq n_{e,k}+b_{e,k}-y_{e,k}\leq K_e .
 \label{eq:storage-bound}
\end{equation}
Paired transfers and route-preserving splits imply
\begin{equation}
 \sum_oQ_{o,k}+\sum_en_{e,k}+Z_k=D_{<k},
 \label{eq:mass-identity}
\end{equation}
also route by route. $Z_k$ is cumulative sink exit; $D_{<k}$ all earlier
exogenous release.

A finite state has a defined unique result: finite source scan, unique
convex allocation vectors, unique capped-sharing vectors and a finite
head-event calculation. Each positive service interval exhausts a
head/quota or reaches step end. Exact ties are handled together; no
mass cutoff or software priority manufactures uniqueness. These are
definition-level arguments, not implementation or traffic-fidelity evidence.

\section{Sensitivities of the update}

Use a fixed structural order and active physical/constraint signature.
Differentiate continuous masses, admissions, queues, histories and head
times. Labels, permissions, access branches and the source-order skeleton
are structural; their switches require branch/one-sided treatment.

For free allocation set ${\cal A}$ and upper-capped set ${\cal U}$,
\begin{align}
 \lambda&=\frac{V-\sum_{\cal U}u_\ell+\sum_{\cal A}h_\ell}
                  {|{\cal A}|},\qquad
 \dot\lambda=\frac{\dot V-\sum_{\cal U}\dot u_\ell+
                                  \sum_{\cal A}\dot h_\ell}
                  {|{\cal A}|},\nonumber\\
 \dot x_\ell&=
 \begin{cases}
 \dot\lambda-\dot h_\ell,&\ell\in{\cal A},\\
 \dot u_\ell,&\ell\in{\cal U},\\
 0,&\text{inactive lower branch}.
 \end{cases}
 \label{eq:allocation-derivative}
\end{align}
On stable fully capped branches $\dot x=\dot u$; stable zero flow has
$\dot x=0$. Differentiate $V$ through its active minimum and $u$ through
its actual residual receiving/access branch. The scan propagates
\begin{equation}
 \dot h\leftarrow\dot h+\dot x,\quad
 \dot R^{\rm rem}\leftarrow\dot R^{\rm rem}-\dot x,\quad
 \dot s\leftarrow\dot s-\dot V,\quad
 \dot q^{\rm rem}=\dot q-\dot V .
 \label{eq:scan-derivative}
\end{equation}
There is no frozen-waiting-vector derivative. A blocked cohort's zero
admission does not erase its retained mass sensitivity. Movement and
physical-prefix switches are not smoothed.

For uncapped leg set ${\cal F}$ and capped set ${\cal C}$,
\begin{equation}
 \dot\theta=
 \frac{\dot W-\sum_{\cal C}\dot C_{ij}
               -\theta\sum_{\cal F}\dot w_{ij}}
      {\sum_{\cal F}w_{ij}},\qquad
 \dot r_{ij}=
 \begin{cases}
 \dot C_{ij},&i\in{\cal C},\\
 w_{ij}\dot\theta+\theta\dot w_{ij},&i\in{\cal F}.
 \end{cases}
 \label{eq:sharing-derivative}
\end{equation}
If the uncapped set is empty, differentiate the cap branches directly;
a stable zero-supply branch has zero grants. No empty-set division is used.
Within-leg sharing uses the unweighted analogue, with the same boundary cases.
For an interior event $\delta=q/r$,
$\dot\delta=\dot q/r-q\dot r/r^2$; the final interval derivative is minus
the preceding duration derivatives. Thus
\begin{equation}
 y=\sum_a r_a\delta_a,\qquad
 \dot y=\sum_a(\dot r_a\delta_a+r_a\dot\delta_a).
 \label{eq:service-derivative}
\end{equation}
Head times cannot be frozen merely because outputs are evaluated at
timestep boundaries.

Fixed delayed queries have fixed interpolation weights but sensitive
history values. For an interior eligibility cut at fixed query $t$,
\begin{equation}
 f=\frac{t-a}{b-a},\quad
 \dot f=\frac{-\dot a-f(\dot b-\dot a)}{b-a},\quad
 \dot q^{\rm ready}=f\dot q+q\dot f.
 \label{eq:eligibility-derivative}
\end{equation}
Use whole/empty branches outside $[a,b]$; point segments need their
boundary branch. For positive entrance $b$,
\begin{equation}
 \dot t_{\rm front}
 =\Delta t\left(\frac{\dot r}{b}-\frac{r\dot b}{b^2}\right),
 \label{eq:entry-derivative}
\end{equation}
with the analogous tail response. No epsilon repairs zero inflow.
Exact-compatible merging must preserve the physical map; a storage-format
change is not automatically a physical kink.

\section{Complete population time and continued state response}

Collect physical stocks, histories, ledgers and sources in $X_k$, with
update $X_{k+1}=F_k(X_k,u_k)$. The inherited cost and its response are
\begin{align}
 J_{\Delta t}
 &=\Delta t\sum_{k=0}^{H-1}
    \left[\sum_o(Q_{o,k}+D_{o,k})+\sum_en_{e,k}\right],
 \label{eq:tst}\\
 P_0&=0,\quad
 P_{k+1}=F_{k,X}P_k+F_{k,u}\dot u_k,\quad
 \dot J=\sum_k(c_{k,X}P_k+c_{k,u}\dot u_k).
 \label{eq:state-sensitivity}
\end{align}
Within control $\dot u_k=d$; in fixed common continuation $\dot u_k=0$,
while $P_k$ keeps evolving. Empty loading/common preparation is independent
of the perturbed candidate; otherwise propagate its response too.
No per-vehicle trajectory gradient is constructed.

At fixed state, the immediate cost partial for current service is zero.
Fork admission also has no direct green input at fixed pre-state: its
feedback comes through earlier queues/history. Stop-line service can
respond directly now. These effects compose into access, spillback,
source waiting and eventual sink exit. Internal transfer alone does not
remove population. For fixed releases,
\begin{equation}
 J_{\Delta t}=\Delta t\sum_{k=0}^{H-1}(D_{\leq k}-Z_k),
 \qquad
 \dot J_{\Delta t}=-\Delta t\sum_{k=0}^{H-1}\dot Z_k .
 \label{eq:exit-objective}
\end{equation}

Use a common integer horizon qualified to clear the compared action
neighborhood, followed by zero occupancy. Finite demand, acyclic routes,
positive continuation service and eventual sink availability support
qualification; no universal numerical horizon or clearance theorem is
asserted. A changing stop counter, residual cutoff or nonclearing case
cannot silently supply complete TST. The derivative domain includes
clearance and unchanged active signatures; switches need directional
branches, not a global smoothness claim.

\section{Matched comparisons and native qualification}

Models share population, geometry/permissions, action, objective,
preparation, continuation and a declared data protocol.
Original fine-grained series may pass through each model's own macro
converter/estimator. Declare retained variables, memory, update schedule
and calibration data. Native subsets are legitimate. Successful macro
inference counts against needing exact order; an exact ordered ledger
retained inside an estimator remains meso information.

Equal current totals do not establish equality of all lawful histories.
State-conditioned replay, observation-updated prediction and forecasting
from a common prefix are distinct tasks. State which observations are
fixed or regenerated for each derivative claim. A later observation
cannot count as a response predicted before it was observed.

Baseline roles remain B0 aggregate, B1 structured movement queues,
B2 responsive lane/storage aggregates and B3 this selective model.
They are not a universally nested information ladder.
Qualify principal instances by soundness/adaptation, task fit and then ease.

B1's working source is
\href{https://arxiv.org/html/2310.12249v2}{Wei et al.'s link queue model},
with moving/queued state and native effective green. Its full causal
short-region/link/node law must be instantiated; the fixed-speed shortcut
is not assumed valid in every spillback regime.

B2's working direction uses
\href{https://www.sciencedirect.com/science/article/pii/S0191261525000578}
{Gong et al.'s lane-exit allocation} with a sourced storage adaptation,
for which \href{https://attap.umd.edu/wp-content/uploads/2022/05/TRC522.pdf}
{Liu--Chang (2011)} is a candidate. Exit allocation is not finite-bay
entry, and native theorems do not automatically qualify that composition.
Exact implementations and adapters still require physical and task-specific
checks; the cited native model is not automatically a qualified baseline.

A V2 comparator failure must be understood before attributing a V3
advantage to retained order.

\section{Definition checks and external usefulness}

Local reference checks cover pool/feeder access, partial admission,
finite storage, current-head
sharing/reuse, route conservation and the source-counted objective.
V0 neutral cases need matched pure service classes and nonbinding or
equivalent receiving; one through/right green does not make differently
constrained exits identical. In that neutral domain the laws reduce to
the same sending/service/receiving minimum. Exact equivalence concerns
a matched aggregate backbone; a source-native propagation difference
must be qualified and cannot be credited to retained order.

Rotating the symmetric scenario and green groups must rotate results.
Equal lane totals do not justify one larger FIFO lane without
corresponding order, readiness, access, service and receiving.
Hand/reference checks establish definition correctness, not realism.

For R2 qualify one fixed executable mapping. With fixed cycle
$C^{\rm cyc}$, lost intervals totaling $L<C^{\rm cyc}$, order and offset,
\begin{equation}
 G_p=(C^{\rm cyc}-L)g_p,\qquad
 s^{\rm eff}_{\ell m}=s^{\rm phys}_{\ell m}
                         (C^{\rm cyc}-L)/C^{\rm cyc}.
 \label{eq:reference-action}
\end{equation}
Declare program admissibility and the common switch to fixed clearance;
do not silently insert minimum green to repair an API.
This matches nominal average opportunity, not microscopic startup/access
or every phase effect. Official
\href{https://sumo.dlr.de/docs/Simulation/Traffic_Lights.html}
{SUMO signal definitions} support the interface, not empirical equivalence.

Use a finite master route population and common desired releases.
Integer reference vehicles are exogenous population, not rounded internal
mass. Fix native lane behavior independently of candidate directions;
do not replay B3's private assignments into only one model.
Use matched predeclared reference seeds and held-out family IDs,
keeping order pairs together.

Reference time is desired-release-to-arrival time:
\texttt{duration} plus \texttt{departDelay}, from official
\href{https://sumo.dlr.de/docs/Simulation/Output/TripInfo.html}{TripInfo}.
Waiting is already inside duration. Retain uninserted/unfinished members
in master-population accounting; clear all demand or report noncompletion,
never a completed-trip-only objective.

For each feasible direction, compare predicted sign/ranking and scale
with actual richer-evaluator complete-time change under the same action
mapping. Disclose observation conditioning. Only after forward/reference
and derivative checks pass should matched optimization and held-out
usefulness be tested.

\section{Scope and evidence limits}

The model specifies eligible packet inspection, movement-specific waiting
priority, actual-admission balancing, physical lane access, elapsed
service, interpolated histories and complete population-time sensitivity.
Its parallel approach lanes retain distinct occupancy, storage and order.
The bay feeder can obstruct one part of the approach while another lane
remains usable, without giving resident traffic an unmodeled bypass.

Feasibility and unique-selection arguments concern the mathematical
operators. The analytic examples illustrate local consequences.
Implementation correctness, numerical derivative agreement, qualified
baselines and external control usefulness require separate evidence.

The source map credits inherited, adapted and newly defined laws.
The LTM basis follows the repository transcription of Yperman (2007) and
UNsim v0.12.0; the ordered route/timing machinery follows its route-mass
specification. Considering independently accessible packets outside the
first sending-budget-sized part of an input list is an explicit urban
adaptation, as is the physical section geometry. No publication is
credited with the whole composed model.
\end{document}
